\chapter*{Introduction generale}
\addcontentsline{toc}{chapter}{Introduction generale}

Le secteur financier connait aujourd'hui une transformation profonde sous l'effet de
la disponibilite massive des donnees, de l'automatisation des traitements et des
progres rapides de l'intelligence artificielle. Les investisseurs disposent d'un
volume important d'informations issues des marches, de l'actualite economique et des
documents reglementaires publies par les entreprises. Cependant, cette abondance ne
garantit pas toujours une meilleure prise de decision : les informations sont souvent
dispersees, difficiles a comparer et parfois marquees par un ecart entre le discours
promotionnel des entreprises et les risques officiellement declares.

Dans ce contexte, les rapports financiers deposes aupres de la SEC, notamment les
rapports annuels 10-K, constituent une source fiable et structuree pour analyser la
situation reelle des societes cotees. Ces documents contiennent des informations
essentielles sur les risques, les performances, les perspectives et la gouvernance des
entreprises. Leur exploitation manuelle reste toutefois complexe en raison de leur
taille, de leur langage technique et du temps necessaire pour identifier les signaux
pertinents. Il devient donc necessaire de concevoir des outils capables de traiter ces
donnees automatiquement et de les transformer en indicateurs exploitables.

Le present projet, intitule \textbf{FinPulse}, s'inscrit dans cette problematique. Il
vise la conception et le developpement d'une plateforme intelligente d'analyse
financiere basee sur les donnees SEC et l'intelligence artificielle. La solution
proposee combine un pipeline d'ingestion de donnees, une base vectorielle, des
mecanismes de recherche augmentee par generation (RAG), un score de coherence appele
\textbf{NCI} et une interface interactive permettant a l'utilisateur d'analyser des
idees d'investissement, de suivre des entreprises et de recevoir des alertes
contextualisees.

L'objectif principal de FinPulse est d'aider l'investisseur a evaluer la coherence
entre une hypothese d'investissement, le discours public d'une entreprise et les
elements reglementaires contenus dans ses rapports officiels. La plateforme permet
ainsi d'interroger un assistant d'analyse, de generer des rapports de strategie,
d'enregistrer des strategies a surveiller et de suivre une watchlist en temps reel
grace a des indicateurs financiers, textuels et comportementaux.

Sur le plan technique, le systeme repose sur une architecture en couches et une
organisation orientee services. L'interface utilisateur est developpee avec Angular,
la logique metier est assuree par un backend Spring Boot, tandis que le traitement des
donnees et des modeles d'intelligence artificielle est pris en charge par un service
Python/FastAPI. La persistance s'appuie sur PostgreSQL et PGVector pour stocker les
donnees structurees ainsi que les embeddings, avec Redis pour l'optimisation des
performances et la gestion du cache.

Ce rapport presente les differentes etapes de realisation du projet. Il introduit
d'abord le contexte general et les objectifs de la solution, puis detaille la
conception fonctionnelle et technique du systeme, notamment les besoins, les
diagrammes UML, l'architecture logicielle et les choix technologiques retenus. Enfin,
il met en evidence l'apport de FinPulse comme outil d'aide a la decision financiere,
en insistant sur la fiabilite des sources, la tracabilite des analyses et la
personnalisation des resultats selon le profil de l'utilisateur.
